\subsection{Частичные пределы. Критерий частичного предела.}

\subsubsection*{Определение}
Число $a$ (или символ $\infty$) называется \textbf{частичным пределом} последовательности $\{a_n\}$, если из неё можно выделить подпоследовательность, сходящуюся к $a$.
\[ a - \text{частичный предел} \Leftrightarrow \exists \{a_{n_k}\}: \lim_{k \to \infty} a_{n_k} = a \]

\textbf{Примеры:}
\begin{itemize}
    \item $a_n = (-1)^n$: Частичные пределы $\{1, -1\}$.
    \item $a_n = n$: Частичный предел $\{+\infty\}$.
    \item $a_n = \sin(\frac{\pi n}{2})$: Частичные пределы $\{0, 1, -1\}$.
\end{itemize}

\subsubsection*{Теорема. Критерий частичного предела}
Число $a$ — частичный предел $\{a_n\} \Leftrightarrow$ в любой окрестности $U_a$ содержится бесконечно много членов последовательности $\{a_n\}$.

\begin{tcolorbox}[title=Доказательство]
	\begin{enumerate}
		\item \textbf{Необходимость ($\Rightarrow$)}
			Пусть $a$ — частичный предел. Тогда существует подпоследовательность $a_{n_k} \to a$.
			По определению предела:
			\[ \forall U_a \;\; \exists K: \; \forall k > K \;\; a_{n_k} \in U_a \]
			Так как $k$ пробегает бесконечное множество значений, то в $U_a$ попадает бесконечно много членов исходной последовательности.

		\item \textbf{Достаточность ($\Leftarrow$)}
			Пусть в любой окрестности точки $a$ бесконечно много членов. Построим подпоследовательность:
			\begin{itemize}
			    \item Возьмем окрестность $(a - 2^{-1}, a + 2^{-1})$. В ней есть элемент $a_{n_1}$.
			    \item Возьмем окрестность $(a - 2^{-2}, a + 2^{-2})$. В ней бесконечно много элементов, выберем $a_{n_2}$ так, чтобы $n_2 > n_1$.
			    \item $\dots$
			    \item На шаге $k$ возьмем окрестность $(a - 2^{-k}, a + 2^{-k})$. Выберем $a_{n_k}$ так, чтобы $n_k > n_{k-1}$.
			\end{itemize}
			Получили подпоследовательность $\{a_{n_k}\}$, для которой $|a_{n_k} - a| < \frac{1}{2^k}$.
			При $k \to \infty$ имеем $a_{n_k} \to a$.
	\end{enumerate}
\end{tcolorbox}

